\documentclass[11pt]{article}
\usepackage[T1]{fontenc}
\usepackage{lmodern,amsmath,amssymb,amsthm,booktabs}
\usepackage[margin=1in]{geometry}
\usepackage[hidelinks]{hyperref}
\newtheorem{proposition}{Proposition}
\newtheorem{corollary}[proposition]{Corollary}
\title{A sharper exponent for integer multiplication}
\author{Douglas Colkitt\thanks{Prepared with assistance from OpenAI Codex. This assistance is not independent review or endorsement by OpenAI.}}
\date{October 6, 2026}
\begin{document}
\maketitle
\begin{abstract}
We give exact certificates for improvements to the parameter choices in
OpenAI's manuscript \emph{Integer multiplication below $n\log n$}.
The strongest supplied substitution uses the same symbolic network construction
at $h=46$, a rational recurrence saving, and a variable stopping exponent to give
$\kappa=29/(5\cdot10^{33})=5.8\cdot10^{-33}$.
This exceeds $2^{-111}$ by a factor of approximately $15.1$ and is more
than $99\%$ of a certified upper bound for the stated parameter family.
We also retain simpler substitutions with earlier interfaces frozen.
This is a conditional consequence of the manuscript's algorithmic
interfaces, not an independent verification of its multiplication theorem.
We also determine the supremum of the symmetric numerical constraint
problem and certify $h=46$ as optimal within the specified network-counting
family when a common recurrence exponent is required.
\end{abstract}

\section{Scope and source}
The source is OpenAI's September 23, 2026 manuscript, released in the
\texttt{openai/math} repository. This note uses commit
\texttt{adc7f1241b42e322a6451854ab7e4b4c146bf78a}.
The repository contains immutable copies of its source and a SHA-256 manifest.
All algorithmic assertions below are conditional on the corresponding
upstream interfaces and their uniform time bounds, on its fixed finite-alphabet
Turing machine with a fixed finite number of one-dimensional tapes. Rational arithmetic
checks prove numerical inequalities; finite tests of selected identities
do not prove the full multitape algorithm.

The starting parameter choices are
\[
 \tau=\sigma=1-2^{-50},\quad \beta=\tfrac12,\quad
 \delta=\tfrac1{16},\quad C_1=20.
\]
The upstream simultaneous-layer proposition permits varying
$c,\epsilon,\lambda,\lambda'$ subject to
\begin{align}
 \max\{\tau,\sigma\}&<\lambda<\lambda'<1,\label{eq:layer1}\\
 \tau(1+2c)&<\lambda,\label{eq:layer2}\\
 \sigma+\tfrac12(1-\sigma)&<\lambda',\qquad 20\epsilon<1.
 \label{eq:layer3}
\end{align}
Its final assembly requires positive $c,\epsilon$, $\epsilon<1/12$,
and positivity of seven margins:
\begin{align*}
 g_1&=1-\epsilon(1+c),&g_2&=\epsilon c(1-\tau),\\
 g_3&=\epsilon(1-\lambda'),&g_4&=1-\tau-\epsilon(2-\tau),\\
 g_5&=\tfrac14-\delta-\tfrac32\epsilon,
 &g_6&=1-\delta-\epsilon,\qquad g_7=\epsilon.
\end{align*}
The final proof reserves half the minimum margin:
$2\kappa\le\min_i g_i$. This absorbs any fixed power of $\log p$ into
$p^\kappa$. This half-margin rule is a convenient convention; mathematically
any fixed $\kappa<\min_i g_i$ would suffice for that absorption.

\section{A substitution with the network frozen}
\begin{proposition}\label{prop:balanced}
Let $0<a\le1/16$ be rational. Suppose the upstream interfaces hold with
$\tau=\sigma=1-a$, $\beta=1/2$, $\delta=1/16$ and $C_1=20$.
Set
\[
 \epsilon=c=a/2,\qquad
 \lambda=1-3a^2/4,\qquad \lambda'=1-a^2/2.
\]
All constraints above hold and $\min_i g_i=a^3/4$. Consequently the
upstream final accounting permits $\kappa=a^3/8$.
\end{proposition}
\begin{proof}
The key layer comparison is
\[
 \tau(1+2c)=(1-a)(1+a)=1-a^2
 <1-3a^2/4<1-a^2/2<1.
\]
Also $1-a<\lambda$ and $1-a/2<\lambda'$.
The dimension and guard constraints follow from
$\epsilon\le1/32<1/12$ and $20\epsilon=10a\le5/8<1$.
Direct substitution gives
\[
\begin{array}{lll}
g_1=1-a/2-a^2/4,&g_2=a^3/4,&g_3=a^3/4,\\[2pt]
g_4=a(1-a)/2,&g_5=3/16-3a/4,&g_6=15/16-a/2,\quad g_7=a/2.
\end{array}
\]
For the stated range, each other margin is greater than $a^3/4$.
\end{proof}

Taking $a=2^{-50}$ gives
\[
 \epsilon=c=2^{-51},\quad
 \lambda=1-3\cdot2^{-102},\quad \lambda'=1-2^{-101},\quad
 \boxed{\kappa=2^{-153}}.
\]
This changes neither finite network nor the exponents in any earlier
proposition. The weaker original proposal,
$\lambda=1-2^{-101}$ and $\lambda'=1-2^{-102}$, instead has
$g_3=2^{-153}$ and permits $\kappa=2^{-154}$ by the same convention.
Both review patches are supplied.

The patch must also change the explicit comparison used to compute
$\lfloor b^\epsilon\rfloor$ from $d^{2^{75}}\le b$ to
$d^{2^{51}}\le b$. Two old displayed estimates cannot be retained:
the intermediate bound $\tau(1+2c)<1-31\cdot2^{-55}$ fails after
substitution, and the new $g_4=2^{-51}-2^{-101}$ is less than
$2^{-51}$. The patches replace the former with its exact identity and
weaken the latter bound to $g_4>2^{-52}$.

\section{Dependencies outside the final inequalities}
The following implications account for uses of the changed exponents
elsewhere in the construction. They use fixed parameters and permit the
common finite cutoff to increase.
\begin{itemize}
\item $\epsilon C_1<1$ implies $d^{C_1}=o(p)$, so the exact intermediate
dyadic formats still have $O(p)$ bits.
\item $\epsilon c>0$ implies $K=\Theta(p^{\epsilon c})$ grows faster than
$\log p$. Every fixed polynomial times $2^{-K}$ tends to zero, including
the repair cost summed over the polynomially many recursive visits.
\item $\epsilon<1$ gives $r=2^{\Theta(p^{1-\epsilon})}$, which dominates
every fixed polynomial in $p$. It pays for per-record and per-line setup.
\item $1-\epsilon-\epsilon c>0$ gives $K=o(\ell)$, as required for the
synthetic-transform layout.
\item $1-2\epsilon>0$ yields the prime-interval hypotheses eventually.
The trial-division setup remains $n^{o(1)}$ since $d\to\infty$.
\item $\epsilon<1/12$ retains
$\alpha=\Theta(p^{1/4+\epsilon/2})<\sqrt p$ and
$\gamma=O(p^{1/2+2\epsilon})=o(p)$.
Thus the upstream error bound $448b\,2^{-b+\gamma}<1/2$ still holds
eventually with $p=6b$.
\item The integer-root exponents are fixed integers. Their size affects
the fixed machine description and constants, not uniformity in input length.
\end{itemize}
No unchanged algorithmic interface is established merely by these
implications. They close the parameter-dependence part of the audit.

\section{The numerical supremum}
\begin{proposition}\label{prop:supremum}
For fixed rational $0<a\le1/32$, with $\tau=\sigma=1-a$,
$\beta=1/2$, $\delta=1/16$ and $C_1=20$, the supremum of
$\min_i g_i$ under the strict layer and assembly constraints is
\[
 G_*(a)=\frac{a^3}{2+a}.
\]
Under the half-margin convention, the supremum of $\kappa$ is
$a^3/(4+2a)$. Neither supremum is attained.
\end{proposition}
\begin{proof}
Put $d_0=1-\lambda'$ and $q=\min\{ac,d_0\}$. The layer conditions give
\[
 d_0< a-2(1-a)c,\qquad d_0<a/2.
\]
The increasing line $ac$ and decreasing line $a-2(1-a)c$ meet at
\[
 c_*=a/(2-a),\qquad q_*=a^2/(2-a).
\]
Strictness implies $q<q_*$. Moreover
\[
 \min_i g_i\le\min\{\epsilon q,\ a-(1+a)\epsilon\}
 \le\frac{a q}{1+a+q}<\frac{a^3}{2+a}.
\]
The middle bound follows by intersecting the two affine functions of
$\epsilon$.

For attainability in the limit, choose rational $0<\xi<1$ and set
\[
 q=(1-\xi)q_*,\quad c=q/a,\quad
 \epsilon=\frac{a}{1+a+q},\quad
 \lambda'=1-q,\quad
 \lambda=\frac{\lambda'+(1-a)(1+2c)}2.
\]
The two endpoints defining $\lambda$ differ by $a\xi>0$.
Also $\lambda>1-a$, $\lambda'<1$ and
$\lambda'>1-a/2$. Since $\epsilon<a\le1/32$, the dimension bound holds
and $20\epsilon<20a\le5/8<1$. The remaining large margins do not bind.
In this family $g_2=g_3=g_4=a q/(1+a+q)$, which tends to $G_*(a)$.
\end{proof}
The certificate uses $\xi=1/100$. Its value of $\kappa$ exceeds
$99/100$ of the half-margin supremum and is strictly below it. Thus
``the strongest safe rational $\kappa$'' is not a maximum in this open
constraint problem. One must specify a slack requirement or a simplicity
criterion for the parameter encoding.

\section{Changing the common recurrence exponent}
Write $v=\binom h3$, $m=h^3$, $N=v^3$, and $I=3v^2$.
The generalized counts from the original construction are
\begin{align*}
 z_b&=3\binom{h-3}{2},&z_c&=\binom{h-3}{3}+3(h-3),\\
 W_b&=2N+I(vz_b+h),&W_c&=2N+I(vz_c+h+1),\\
 L_b&=Ih^2,&L_c&=Ih(h+1),\\
 \eta_b&=\frac{N-2L_b}{W_bm},&
 \eta_c&=\frac{2N-2L_c}{W_cm}.
\end{align*}
If $\log m<L$ and both deficits exceed $La$, then
\[
 m^{1-a}=m e^{-a\log m}>m(1-aL)>m(1-\eta),
\]
which supplies the strict recurrence comparison for either network.
All comparisons used here have rational certificates. A positive Taylor
partial sum exceeding $m$ proves $e^L>m$ without floating-point arithmetic.

For $h=100$, the published deficits exceed $14\cdot2^{-42}$ and
$\log(10^6)<14$. Thus the same network supports the comparison with
$a=2^{-42}$. Proposition~\ref{prop:balanced} then gives the candidate
$\kappa=2^{-129}$, improving on the unbalanced $2^{-130}$ candidate.
The separate \texttt{same-network-129.patch} updates the proposition statements
that fix the recurrence exponents and the displayed exponent comparison.

For $h=46$, the formulas give
\[
 m=97336,\quad \eta_b=\frac9{43518487588},\quad
 \eta_c=\frac7{22253827054}.
\]
Both deficits exceed $12\cdot2^{-36}$ and $\log m<12$.
The guard calculation has
$B=s_c+64(W_c+m+1)^3$ with $m^{10}<B<m^{11}$, so $C_1=20$
still suffices. Proposition~\ref{prop:balanced} gives the numerical
candidate $\kappa=2^{-111}$. The \texttt{h46-111.patch} updates six source
files with these counts, recurrence exponents, and assembly parameters.

At this $h$, the rational form $I-J/9$ has exceptional eigenvalue
$1-h/9=-37/9\ne0$. Triple vectors retain rational norm two and
binary norm one. The support arguments for nonalternating residuals
still have spare coordinates because $h>6$. The tensor depth remains
three and the endpoint weight remains $27$. These observations explain
why the symbolic construction is a plausible extension; the counting
certificate alone does not certify that extension or its implementation.

\section{A certified discrete optimum}
For the stated admissible family ($h\ge7$, $h\ne9$, and both
$L_b,L_c<N/2$), the supremum of the shared recurrence saving is
\[
 A(h)=\frac{-\log(1-\min(\eta_b,\eta_c))}{\log(h^3)}.
\]
The supplied search certifies $A(46)>A(h)$ for all other admissible
integers $h$. It compares rational logarithm enclosures for $h<200$,
and bounds the infinite tail as follows:
\[
 \eta_b<\frac1{3z_bm},\qquad
 A(h)\le\frac{-\log(1-\eta_b)}{\log m}
 <\frac1{3z_bm-1}\quad(h\ge200).
\]
Both $z_b$ and $m$ increase with $h$, and the bound at $200$ is already
below the certified lower bound for $A(46)$.

For $1\le x\le2$, the logarithm enclosures use
\[
 z=\frac{x-1}{x+1},\quad
 2\sum_{j=0}^{J-1}\frac{z^{2j+1}}{2j+1}
 \le\log x\le
 2\sum_{j=0}^{J-1}\frac{z^{2j+1}}{2j+1}
 +\frac{2z^{2J+1}}{(2J+1)(1-z^2)}.
\]
Integer logarithms are reduced using
$\log m=k\log2+\log(m/2^k)$.
This proves an optimum for a common recurrence exponent within these
counting formulas. It does not optimize separate $\tau,\sigma$, altered
tensor depth, different networks, or the full space of multiplication algorithms.

\section{Rational parameters and variable stopping exponent}
The sharper rational choice
\[
 a=\frac9{500000000000},\qquad L=\frac{5743}{500}
\]
is certified at $h=46$: both exact rank deficits exceed $aL$, and the
logarithm enclosure above with $J=24$ proves $\log(97336)<L$.
Hence both network comparisons support $\tau=\sigma=1-a$.

\begin{proposition}[Variable stopping exponent]
Conditional on the same network and stream interfaces, the proof of the
simultaneous-layer proposition works for any fixed rational $0<\beta<1$
provided its displayed layer inequalities hold with that $\beta$.
\end{proposition}
\begin{proof}
At an internal node $e\ge d^\beta$, the packed overhead satisfies
$K\le d^c\le e^{c/\beta}$ and hence is $O(e^\lambda)$ when
$\tau(1+c/\beta)<\lambda$. The internal geometric sum uses
$s/W\le m^\sigma<m^\lambda$. At a leaf $u<d^\beta$, the normalized
cost summed over leaves is at most
$O(e^\sigma u^{1-\sigma})\le O(e^\sigma d^{\beta(1-\sigma)})$.
The logarithmic number of base-$m$ pieces is absorbed when
$\lambda'>\max\{\lambda,\sigma+\beta(1-\sigma)\}$.

The initial row preprocessing remains $O(V\log d)$. The depth is at most
$\lceil\log_m d\rceil$ and each leaf has $e\le d$, so the same row
padding and guard-depth bounds apply. For $\beta=u/v$ the stopping test
is the exact integer comparison $e^v<d^u$. The powers are fixed and their
polynomial descriptor work is absorbed by the length of each polynomial record
$r=2^{\Theta(p^{1-\epsilon})}$ as before. No tape-count change is needed.
\end{proof}

\begin{proposition}[Strongest supplied witness]
With the rational $a$ above, set
\begin{gather*}
 \tau=\sigma=1-a,\quad \delta=\tfrac1{16},\quad C_1=20,\\
 \beta=\frac{999}{1000},\quad
 \epsilon=\frac{999a}{1000},\quad c=\frac{998a}{1000},\\
 \lambda'=1-ac,\qquad
 \lambda=\frac{\lambda'+(1-a)(1+c/\beta)}2.
\end{gather*}
All layer, assembly, and recorded side conditions hold, and the
conditional final bound permits
\[
 \boxed{\kappa=\frac{29}{5\cdot10^{33}}=5.8\cdot10^{-33}}.
\]
\end{proposition}
\begin{proof}
The endpoint gap defining $\lambda$ is
\[
 \lambda'-(1-a)(1+c/\beta)
 =\frac{a}{999}+\frac{998a^2}{999000}>0.
\]
Both endpoints exceed $1-a$. The leaf inequality is
$1-a/1000<1-998a^2/1000$, which holds since $998a<1$.
The dimension and guard requirements hold because $\epsilon<a<1/20$.
Direct substitution gives
\[
 G:=g_2=g_3=\frac{498501}{500000}a^3
   =5.814515664\cdot10^{-33}>\kappa.
\]
Here $g_1,g_6>1/2$, $g_5>1/8$, $g_4>a/2000$ and $g_7>a/2$;
all exceed $G$. The other side conditions in the dependency section hold
by substitution. Put $\rho=G-\kappa>0$. Every fixed logarithmic factor
is eventually at most $p^\rho$, giving
$p^{1-G}(\log p)^{O(1)}=O(p^{1-\kappa})$.
\end{proof}

Two simpler intermediate witnesses are also supplied:
\[
\begin{array}{c|c|c|c|c|c}
\kappa&\beta&\epsilon&c&\lambda&\lambda'\\\hline
2^{-109}&1/2&3a/4&a/2&1-3a^2/4&1-a^2/2\\
2^{-108}&3/4&3a/4&3a/4&1-7a^2/8&1-3a^2/4
\end{array}
\]
Their minimum margins are respectively $3a^3/8$ and $9a^3/16$,
strictly above the listed $\kappa$. Each patch updates the rational
integer-power setup, network comparison, and final absorption argument.

\section{A ceiling for further parameter tuning}
This ceiling allows variable $\beta$ and unequal $\tau,\sigma$, but retains
the stated network rank bounds and seven-margin accounting.
Write $a=1-\tau$. Since $g_4>0$, we have $\epsilon<a/(1+a)<a$.
If $c<a$ then $ac<a^2$. If $c\ge a$, the packed-overhead condition and
$\beta<1$ give
\[
 1-\lambda'<1-(1-a)(1+c/\beta)
 <1-(1-a)(1+a)=a^2.
\]
In either case
\[
 \kappa<\min\{g_2,g_3\}
 =\epsilon\min\{ac,1-\lambda'\}<a^3.
\]
The certified search was repeated for the bit-network saving alone,
\[
 A_b(h)=\frac{-\log(1-\eta_b(h))}{\log(h^3)}.
\]
Its unique maximum is again at $h=46$ over the stated admissible integers;
the same decreasing tail bound covers $h\ge200$. The cubed upper
logarithm enclosure certifies
\[
 \kappa<A_b(46)^3<5.838\cdot10^{-33}<2^{-107}.
\]
The supplied $5.8\cdot10^{-33}$ exceeds $99\%$ of that upper enclosure.
Thus a further improvement to $2^{-107}$ requires changing these rank
bounds, the admissible network construction, or the cost accounting.
This is not a lower bound for arbitrary multiplication algorithms.

\section{Reproduction and remaining proof work}
Run \texttt{make verify} to regenerate exact certificates, all seven review
patches, and the finite checks. Run \texttt{make note} to build this note.
The dependency audit is recorded in \texttt{docs/audit.md}.

The first patch changes three source files and no substantive lemma.
The arithmetic and parameter implications are reviewable independently.
Establishing the unconditional multiplication result still requires accepting
or independently validating the upstream finite-network interfaces,
fixed-tape implementations, and analytic reductions. No complete Lean
formalization or executable implementation of the full multiplication
machine is supplied here. No priority or publication claim is made.

\begin{thebibliography}{9}
\bibitem{openai} OpenAI, \emph{Integer multiplication below $n\log n$},
September 23, 2026. Source in
\url{https://github.com/openai/math/tree/adc7f1241b42e322a6451854ab7e4b4c146bf78a}.
\bibitem{hh} D. Harvey and J. van der Hoeven,
\emph{Integer multiplication in time $O(n\log n)$},
Annals of Mathematics 193 (2021), 563--617.
\url{https://www.texmacs.org/joris/nlogn/nlogn.pdf}.
\end{thebibliography}
\end{document}
